\documentclass[11pt]{article}

\usepackage{amsmath,amssymb,amsthm,mathtools}
\usepackage[margin=1in]{geometry}

\newtheorem{theorem}{Theorem}[section]
\newtheorem{lemma}[theorem]{Lemma}
\newtheorem{corollary}[theorem]{Corollary}
\newtheorem{remark}[theorem]{Remark}
\newtheorem{definition}[theorem]{Definition}

\newcommand{\dd}{\,d}

\begin{document}

\section{Triangle books with leaves rooted at page vertices}

For a finite graph $H$ and $0\leq p<1$, put
\[
 \Phi_H(p)
 :=(1-p)^{v(H)}
 \chi_H\!\left(\frac{1}{1-p}\right),
\]
and use its polynomial continuation at $p=1$.  We prove a positive
theorem for triangle books whose pendant leaves are distributed among
the page vertices.  This is complementary to the previously proved
two-root theorem, in which leaves are attached to the two ends of the
common spine edge.

Throughout, $(\Omega,\mathcal A,\mu)$ is an arbitrary probability
space and $W:\Omega^2\to[0,1]$ is a symmetric jointly measurable
graphon.  Write
\[
 p=t(K_2,W)=\int_{\Omega^2}W(x,y)\dd\mu(x)\dd\mu(y),
 \qquad
 d(x)=\int_\Omega W(x,y)\dd\mu(y).
\]


%%%%%%%%%%%%%%%%%%%%%%%%%%%%%%%%%%%%%%%%%%%%%%%%%%%%%%%%%%%%%%%%%%%%%%%%%%%%%%%
\section{The graph family and its exact target}
%%%%%%%%%%%%%%%%%%%%%%%%%%%%%%%%%%%%%%%%%%%%%%%%%%%%%%%%%%%%%%%%%%%%%%%%%%%%%%%

\begin{definition}[Page-rooted leaf book]\label{def:family}
Let $m\geq1$ and let $\mathbf{k}=(k_1,\ldots,k_m)$ be a vector of
nonnegative integers.  The graph $P_{m;\mathbf{k}}$ has two adjacent
\emph{spine vertices} $u,v$, page vertices $z_1,\ldots,z_m$, and
$k_i$ private pendant leaves adjacent only to $z_i$.  Each page vertex
is adjacent to both spine vertices, and there are no other edges.

Put
\[
 n:=k_1+\cdots+k_m.
\]
Then
\[
 v(P_{m;\mathbf{k}})=m+n+2,
 \qquad
 e(P_{m;\mathbf{k}})=2m+n+1.
\]
\end{definition}

\begin{lemma}[Chromatic data]\label{lem:chromatic}
For every $m\geq1$ and every $\mathbf{k}\in\mathbb N^m$,
\[
 \chi(P_{m;\mathbf{k}})=3
\]
and
\begin{equation}
 \chi_{P_{m;\mathbf{k}}}(x)
 =x(x-1)^{n+1}(x-2)^m.
 \label{eq:chromatic}
\end{equation}
Consequently, for every $p\in[0,1]$,
\begin{equation}
 \Phi_{P_{m;\mathbf{k}}}(p)
 =p^{n+1}(2p-1)^m.
 \label{eq:target}
\end{equation}
\end{lemma}

\begin{proof}
Choose the two distinct colours on the ordered spine in $x(x-1)$
ways.  Each page vertex then has $x-2$ choices, independently of the
other pages.  After a page colour has been chosen, each of its private
leaves has $x-1$ choices.  This proves \eqref{eq:chromatic}.  The graph
contains a triangle and the same construction gives a $3$-colouring,
so its chromatic number is $3$.

For $q=1-p>0$, substitution gives
\begin{align*}
 \Phi_{P_{m;\mathbf{k}}}(p)
 &=q^{m+n+2}\frac1q
   \left(\frac pq\right)^{n+1}
   \left(\frac{2p-1}{q}\right)^m
 \\
 &=p^{n+1}(2p-1)^m.
\end{align*}
The last expression is a polynomial in $p$ and therefore also gives
the continued value $1$ at $p=1$.
\end{proof}


%%%%%%%%%%%%%%%%%%%%%%%%%%%%%%%%%%%%%%%%%%%%%%%%%%%%%%%%%%%%%%%%%%%%%%%%%%%%%%%
\section{A weighted rooted-triangle inequality for real exponents}
%%%%%%%%%%%%%%%%%%%%%%%%%%%%%%%%%%%%%%%%%%%%%%%%%%%%%%%%%%%%%%%%%%%%%%%%%%%%%%%

Define
\[
 \tau(x):=\int_{\Omega^2}
 W(x,y)W(x,z)W(y,z)\dd\mu(y)\dd\mu(z)
\]
and, for every real $s\geq0$,
\[
 M_s:=\int_\Omega d(x)^s\dd\mu(x),
\]
with the convention $d(x)^0=1$, including when $d(x)=0$.

\begin{lemma}[Weighted rooted triangle]\label{lem:weighted-triangle}
If $p>0$, then for every real $h\geq0$,
\begin{equation}
 \int_\Omega d(x)^h\tau(x)\dd\mu(x)
 \geq
 \frac{2p-1}{p}M_{h+2}.
 \label{eq:weighted-triangle}
\end{equation}
\end{lemma}

\begin{proof}
For a bounded measurable kernel $K$, write
\[
 (T_Kf)(x):=\int_\Omega K(x,y)f(y)\dd\mu(y),
 \qquad
 \langle f,g\rangle:=\int_\Omega f(x)g(x)\dd\mu(x).
\]
Set
\[
 I_h:=\int d^h\tau,
 \qquad
 J_h:=\langle d^h,T_Wd\rangle.
\]
The elementary inequality $ab\geq a+b-1$ on $[0,1]^2$, applied to
$W(x,y)$ and $W(x,z)$ and then multiplied by the nonnegative factor
$d(x)^hW(y,z)$, gives after Tonelli and symmetry
\begin{equation}
 I_h\geq2J_h-pM_h.
 \label{eq:first-reduction}
\end{equation}

Put $U=1-W$, $u=T_U\mathbf1=1-d$, and $q=1-p$.  Since
\[
 T_Wd=d-q\mathbf1+T_Uu,
\]
we have
\begin{equation}
 J_h=M_{h+1}-qM_h+\langle d^h,T_Uu\rangle.
 \label{eq:J-expansion}
\end{equation}
Symmetry of $U$ gives the exact identity
\begin{align*}
 &2\left(\langle d^h,T_Uu\rangle
       -\int_\Omega d(x)^hu(x)^2\dd\mu(x)\right)
 \\
 &\quad=
 \int_{\Omega^2}U(x,y)
 \bigl(u(y)-u(x)\bigr)
 \bigl(d(x)^h-d(y)^h\bigr)
 \dd\mu(x)\dd\mu(y).
\end{align*}
Because $d=1-u$ and $s\mapsto s^h$ is nondecreasing on $[0,1]$,
the integrand on the right is nonnegative.  Therefore
\[
 \langle d^h,T_Uu\rangle
 \geq\int d^h(1-d)^2
 =M_h-2M_{h+1}+M_{h+2}.
\]
Substituting this into \eqref{eq:J-expansion}, and then into
\eqref{eq:first-reduction}, yields
\[
 I_h\geq2M_{h+2}-2M_{h+1}+pM_h.
\]
Finally,
\begin{align*}
 &2M_{h+2}-2M_{h+1}+pM_h
 -\frac{2p-1}{p}M_{h+2}
 \\
 &\quad=
 \frac1p\int_\Omega d(x)^h(d(x)-p)^2\dd\mu(x)
 \geq0.
\end{align*}
This proves \eqref{eq:weighted-triangle}.  Notice that no integrality
of $h$ was used; this is the point needed below.
\end{proof}


%%%%%%%%%%%%%%%%%%%%%%%%%%%%%%%%%%%%%%%%%%%%%%%%%%%%%%%%%%%%%%%%%%%%%%%%%%%%%%%
\section{Compression of the page-rooted moments}
%%%%%%%%%%%%%%%%%%%%%%%%%%%%%%%%%%%%%%%%%%%%%%%%%%%%%%%%%%%%%%%%%%%%%%%%%%%%%%%

For $x,y\in\Omega$ and real $s\geq0$, define
\[
 H_s(x,y):=
 \int_\Omega W(x,z)W(y,z)d(z)^s\dd\mu(z).
\]

\begin{lemma}[H\"older compression]\label{lem:holder-compression}
Let $k_1,\ldots,k_m\geq0$ be real numbers and put
\[
 \alpha:=\frac{k_1+\cdots+k_m}{m}.
\]
Then, pointwise for all $x,y\in\Omega$,
\begin{equation}
 \prod_{i=1}^m H_{k_i}(x,y)\geq H_\alpha(x,y)^m.
 \label{eq:holder-compression}
\end{equation}
\end{lemma}

\begin{proof}
Apply generalized H\"older inequality to the $m$ functions
\[
 z\longmapsto
 \bigl(W(x,z)W(y,z)d(z)^{k_i}\bigr)^{1/m}.
\]
Their pointwise product is
$W(x,z)W(y,z)d(z)^\alpha$.  Thus
\[
 H_\alpha(x,y)
 \leq\prod_{i=1}^m H_{k_i}(x,y)^{1/m},
\]
which is equivalent to \eqref{eq:holder-compression}.  The convention
$d^0=1$ handles zero exponents, and no quotient is used.
\end{proof}


%%%%%%%%%%%%%%%%%%%%%%%%%%%%%%%%%%%%%%%%%%%%%%%%%%%%%%%%%%%%%%%%%%%%%%%%%%%%%%%
\section{The page-rooted leaf theorem}
%%%%%%%%%%%%%%%%%%%%%%%%%%%%%%%%%%%%%%%%%%%%%%%%%%%%%%%%%%%%%%%%%%%%%%%%%%%%%%%

\begin{theorem}[Leaves on the pages of a triangle book]
\label{thm:main}
Let $m\geq1$ and $\mathbf{k}=(k_1,\ldots,k_m)\in\mathbb N^m$.
For every graphon $W$ of edge density $p\in[1/2,1]$,
\[
 \boxed{
 t(P_{m;\mathbf{k}},W)
 \geq p^{n+1}(2p-1)^m
 =\Phi_{P_{m;\mathbf{k}}}(p).
 }
\]
This is the full interval required by
$\chi(P_{m;\mathbf{k}})=3$.
\end{theorem}

\begin{proof}
Assume first that $1/2\leq p<1$; in particular $p>0$.  Integrating
out every private leaf and then every page vertex gives the exact
identity
\begin{equation}
 t(P_{m;\mathbf{k}},W)
 =\int_{\Omega^2}W(x,y)
   \prod_{i=1}^m H_{k_i}(x,y)
   \dd\mu(x)\dd\mu(y).
 \label{eq:density-identity}
\end{equation}
All integrands are nonnegative and bounded, so Tonelli justifies this
factorization on an arbitrary probability space.

Put $\alpha=n/m$.  By Lemma~\ref{lem:holder-compression},
\[
 t(P_{m;\mathbf{k}},W)
 \geq\int_{\Omega^2}W(x,y)H_\alpha(x,y)^m
 \dd\mu(x)\dd\mu(y).
\]
Use the edge-biased probability measure
\[
 \dd\rho(x,y):=\frac{W(x,y)}p
 \dd\mu(x)\dd\mu(y).
\]
Jensen's inequality for the convex function $s\mapsto s^m$ gives
\begin{align}
 t(P_{m;\mathbf{k}},W)
 &\geq p\bigl(\mathbb E_\rho H_\alpha\bigr)^m
 \notag\\
 &=p^{1-m}
 \left(\int_{\Omega^2}W(x,y)H_\alpha(x,y)
 \dd\mu(x)\dd\mu(y)\right)^m.
 \label{eq:edge-jensen}
\end{align}
Expanding $H_\alpha$ and applying Tonelli once more gives
\[
 \int_{\Omega^2}W(x,y)H_\alpha(x,y)\dd\mu(x)\dd\mu(y)
 =\int_\Omega d(z)^\alpha\tau(z)\dd\mu(z).
\]
Lemma~\ref{lem:weighted-triangle} at the real exponent $h=\alpha$
and Jensen's inequality for $s\mapsto s^{\alpha+2}$ imply
\[
 \int d^\alpha\tau
 \geq\frac{2p-1}{p}M_{\alpha+2}
 \geq\frac{2p-1}{p}p^{\alpha+2}
 =(2p-1)p^{\alpha+1}.
\]
Substitution into \eqref{eq:edge-jensen}, using $m\alpha=n$, yields
\begin{align*}
 t(P_{m;\mathbf{k}},W)
 &\geq p^{1-m}
 \bigl((2p-1)p^{\alpha+1}\bigr)^m
 \\
 &=p^{n+1}(2p-1)^m.
\end{align*}
By Lemma~\ref{lem:chromatic}, this is exactly the required target.

If $p=1$, then $0\leq W\leq1$ and $\int W=1$ imply
$W=1$ almost everywhere.  Every finite homomorphism density is then
$1$, equal to the polynomially continued target.
\end{proof}


%%%%%%%%%%%%%%%%%%%%%%%%%%%%%%%%%%%%%%%%%%%%%%%%%%%%%%%%%%%%%%%%%%%%%%%%%%%%%%%
\section{A conditional cone consequence}
%%%%%%%%%%%%%%%%%%%%%%%%%%%%%%%%%%%%%%%%%%%%%%%%%%%%%%%%%%%%%%%%%%%%%%%%%%%%%%%

For $d(x)>0$, define the probability measure
\[
 \dd\mu_x(y):=\frac{W(x,y)}{d(x)}\dd\mu(y),
\]
and regard the same kernel $W$ as a graphon $W_x$ on
$(\Omega,\mathcal A,\mu_x)$.  Its edge density is
\[
 z_x=t(K_2,W_x)=\frac{\tau(x)}{d(x)^2}.
\]
Put $z_x=0$ when $d(x)=0$.

\begin{lemma}[Conditional cone lemma]\label{lem:cone}
Let $F$ be a graph on $N\geq2$ vertices and put
$\phi(z)=\Phi_F(z)$.  Suppose that, for some $a\in[0,1)$:
\begin{enumerate}
 \item every graphon $V$ on every probability space, of edge density
 $z\geq a$, satisfies $t(F,V)\geq\phi(z)$;
 \item $\phi\geq0$, $\phi'\geq0$, and $\phi''\geq0$ on $[a,1]$;
 \item $\phi(a)=0$.
\end{enumerate}
Then every graphon $W$ of edge density
\[
 p\geq\frac1{2-a}
\]
satisfies
\[
 t(K_1\vee F,W)\geq\Phi_{K_1\vee F}(p).
\]
\end{lemma}

\begin{proof}
The endpoint $p=1$ is immediate, so assume $p<1$ and put
$c=2-1/p$.  The density hypothesis is exactly $c\geq a$, while
$c\leq1$.  Every graphon $V$ of arbitrary edge density $z\in[0,1]$
satisfies the global tangent bound
\begin{equation}
 t(F,V)\geq\phi(c)+\phi'(c)(z-c).
 \label{eq:global-tangent}
\end{equation}
Indeed, for $z\geq a$ this follows from the assumed bound and
convexity.  For $z<a$, the tangent line is nonpositive at $a$ and has
nonnegative slope, so it remains nonpositive; then
$t(F,V)\geq0$ proves \eqref{eq:global-tangent}.

Conditioning on the image of the cone vertex gives
\[
 t(K_1\vee F,W)=\int_\Omega d(x)^N t(F,W_x)\dd\mu(x).
\]
The set $\{d=0\}$ contributes zero and entails no division.  Applying
\eqref{eq:global-tangent} in each nonzero link yields
\begin{align*}
 t(K_1\vee F,W)
 &\geq M_N\phi(c)
 +\phi'(c)\left(\int d(x)^Nz_x\dd\mu(x)-cM_N\right).
\end{align*}
By Lemma~\ref{lem:weighted-triangle} with $h=N-2$,
\[
 \int d(x)^Nz_x\dd\mu(x)
 =\int d(x)^{N-2}\tau(x)\dd\mu(x)
 \geq cM_N.
\]
Thus the second term is nonnegative.  Since $\phi(c)\geq0$ and
$M_N\geq p^N$ by Jensen,
\[
 t(K_1\vee F,W)\geq p^N\phi(c).
\]

Let $q=1-p>0$.  Since $1-c=q/p$ and
$\chi_{K_1\vee F}(x)=x\chi_F(x-1)$,
\[
 p^N\phi(c)
 =q^N\chi_F\!\left(\frac pq\right)
 =q^{N+1}\chi_{K_1\vee F}\!\left(\frac1q\right)
 =\Phi_{K_1\vee F}(p).
\]
\end{proof}

\begin{theorem}[Cones over page-rooted leaf books]\label{thm:cones}
For every $m\geq1$ and $\mathbf{k}\in\mathbb N^m$, the cone
\[
 C_{m;\mathbf{k}}:=K_1\vee P_{m;\mathbf{k}}
\]
satisfies, for every graphon of edge density $p\in[2/3,1]$,
\[
 \boxed{
 t(C_{m;\mathbf{k}},W)
 \geq p(2p-1)^{n+1}(3p-2)^m
 =\Phi_{C_{m;\mathbf{k}}}(p).
 }
\]
Moreover,
\[
 \chi_{C_{m;\mathbf{k}}}(x)
 =x(x-1)(x-2)^{n+1}(x-3)^m,
 \qquad
 \chi(C_{m;\mathbf{k}})=4.
\]
Thus $[2/3,1]$ is the complete required interval.
\end{theorem}

\begin{proof}
The chromatic identity follows from
$\chi_{K_1\vee F}(x)=x\chi_F(x-1)$ and
Lemma~\ref{lem:chromatic}.  Direct substitution into the definition
of $\Phi$ gives the displayed target.

Theorem~\ref{thm:main} supplies the base bound for
\[
 \phi(z)=z^{n+1}(2z-1)^m
 \qquad (z\in[1/2,1]).
\]
On this interval, $\phi$ is a product of nonnegative nondecreasing
affine factors.  Its first and second derivatives are sums of products
of nonnegative factors and positive slopes.  Hence
$\phi,\phi',\phi''\geq0$, and $\phi(1/2)=0$.  Apply
Lemma~\ref{lem:cone} with $a=1/2$; its threshold is
\[
 \frac1{2-1/2}=\frac23.
\]
\end{proof}


%%%%%%%%%%%%%%%%%%%%%%%%%%%%%%%%%%%%%%%%%%%%%%%%%%%%%%%%%%%%%%%%%%%%%%%%%%%%%%%
\section{Thresholds and equality}
%%%%%%%%%%%%%%%%%%%%%%%%%%%%%%%%%%%%%%%%%%%%%%%%%%%%%%%%%%%%%%%%%%%%%%%%%%%%%%%

At the required threshold $p=1/2$, the target is zero.  The proof
above is still valid because the only normalized measure divides by
$p$, not by $2p-1$ or by a triangle density.  Equivalently, the
threshold conclusion follows immediately from nonnegativity.

For every integer $\ell\geq2$, the balanced complete $\ell$-partite
graphon $T_\ell$ has $p=1-1/\ell$ and
\[
 t(P_{m;\mathbf{k}},T_\ell)
 =\frac{\chi_{P_{m;\mathbf{k}}}(\ell)}
        {\ell^{m+n+2}}
 =\Phi_{P_{m;\mathbf{k}}}(p).
\]
Thus the theorem has the expected Tur\'an equality cases, including
$T_2$ at the threshold.  At $p=1$, both sides are again $1$.

For the cone family, the target vanishes at $p=2/3$, and the balanced
complete tripartite graphon attains equality because every cone contains
a $K_4$.  More generally, the same colouring-density calculation gives
equality on every balanced multipartite grid point $p=1-1/\ell$ with
$\ell\geq3$.


%%%%%%%%%%%%%%%%%%%%%%%%%%%%%%%%%%%%%%%%%%%%%%%%%%%%%%%%%%%%%%%%%%%%%%%%%%%%%%%
\section{Catalogue consequences}
%%%%%%%%%%%%%%%%%%%%%%%%%%%%%%%%%%%%%%%%%%%%%%%%%%%%%%%%%%%%%%%%%%%%%%%%%%%%%%%

\begin{corollary}[New positive Atlas cases]\label{cor:atlas}
The following previously open graphs satisfy
$t(H,W)\geq\Phi_H(p)$ for every graphon throughout the full required
interval $p\in[1/2,1]$:
\[
\begin{array}{c|c|c|c|c|c|c}
\text{Atlas}&\text{graph6}&v&e&\mathbf{k}&\chi_H(x)&\Phi_H(p)\\ \hline
41&\texttt{Db[}&5&6&(1,0)&x(x-1)^2(x-2)^2&p^2(2p-1)^2\\
114&\texttt{EJwG}&6&7&(2,0)&x(x-1)^3(x-2)^2&p^3(2p-1)^2\\
138&\texttt{Ez`\_}&6&8&(1,0,0)&x(x-1)^2(x-2)^3&p^2(2p-1)^3
\end{array}
\]
For Atlas 41 and 114 one has $m=2$; for Atlas 138 one has $m=3$.
\end{corollary}

\begin{proof}
For an unambiguous finite identification, use the canonical labelling
from Definition~\ref{def:family}: the spine is $01$, the pages are
$2,3,\ldots,m+1$, and the leaves follow the page vertices.

For Atlas 41 the canonical edge list is
\[
 01,02,03,12,13,24.
\]
It is carried to the NetworkX Atlas edge list
\[
 01,13,14,23,24,34
\]
by
\[
 0\mapsto3,\quad1\mapsto4,\quad2\mapsto1,\quad
 3\mapsto2,\quad4\mapsto0.
\]

For Atlas 114 the canonical edge list is
\[
 01,02,03,12,13,24,25.
\]
It is carried to
\[
 04,12,13,14,23,24,45
\]
by
\[
 0\mapsto1,\quad1\mapsto2,\quad2\mapsto4,\quad
 3\mapsto3,\quad4\mapsto0,\quad5\mapsto5.
\]

For Atlas 138 the canonical edge list is
\[
 01,02,03,04,12,13,14,25.
\]
It is carried to
\[
 01,02,04,12,13,15,23,25
\]
by
\[
 0\mapsto1,\quad1\mapsto2,\quad2\mapsto0,\quad
 3\mapsto3,\quad4\mapsto5,\quad5\mapsto4.
\]
These relabellings verify the structural hypotheses exactly.  The
orders, sizes, chromatic polynomials, and targets now follow from
Lemma~\ref{lem:chromatic}, and Theorem~\ref{thm:main} proves all three
rows.
\end{proof}

\begin{corollary}[New positive cone: Atlas 193]\label{cor:atlas-cone}
NetworkX Atlas graph $193$, with graph6 code \texttt{En\char125 G},
satisfies, for every graphon of edge density $p\in[2/3,1]$,
\[
 t(F_{193},W)
 \geq p(2p-1)^2(3p-2)^2
 =\Phi_{F_{193}}(p).
\]
It has order $6$, size $11$, chromatic number $4$, and
\[
 \chi_{F_{193}}(x)=x(x-1)(x-2)^2(x-3)^2.
\]
\end{corollary}

\begin{proof}
This graph is the cone over
$P_{2;(1,0)}$, the representative used for Atlas 41 above.  Add cone
vertex $5$ to its canonical labelling.  The resulting edge list is
\[
 01,02,03,05,12,13,15,24,25,35,45.
\]
The NetworkX Atlas 193 edge list is
\[
 01,03,04,05,12,13,14,23,24,34,45.
\]
The first is carried to the second by
\[
 0\mapsto1,\quad1\mapsto3,\quad2\mapsto0,\quad
 3\mapsto2,\quad4\mapsto5,\quad5\mapsto4.
\]
Theorem~\ref{thm:cones} with $(m,\mathbf{k})=(2,(1,0))$ now gives
the claimed polynomial, target, chromatic number, and graphon bound.
\end{proof}

The same family also recovers Atlas 15, 34, and 92 when $m=1$, and
Atlas 115 when $(m,\mathbf{k})=(2,(1,1))$.  Those rows were already
positive by, respectively, the rooted triangle--tree and edge
self-amalgamation theorems.  Exact enumeration of all members of the
base and cone families with at most six vertices shows that the three
rows in Corollary~\ref{cor:atlas} and the cone in
Corollary~\ref{cor:atlas-cone} are the only new catalogue consequences.


%%%%%%%%%%%%%%%%%%%%%%%%%%%%%%%%%%%%%%%%%%%%%%%%%%%%%%%%%%%%%%%%%%%%%%%%%%%%%%%
\section{Accepted inputs, hidden-assumption audit, and limitations}
%%%%%%%%%%%%%%%%%%%%%%%%%%%%%%%%%%%%%%%%%%%%%%%%%%%%%%%%%%%%%%%%%%%%%%%%%%%%%%%

The proof is self-contained at the graphon level.  Its only analytic
inputs are the elementary inequality $ab\geq a+b-1$, generalized
H\"older, Jensen, Tonelli--Fubini, and monotonicity of real powers on
$[0,1]$.  The weighted rooted-triangle estimate is reproved for every
real exponent $h\geq0$, rather than cited only in its previously used
integer form.

No atomlessness, standard-Borel representation, regularity,
finite-step approximation, injective-copy interpretation, minimum
degree, or pointwise lower bound on $d$ is assumed.  All integrands are
measurable and bounded.  Zero degrees cause no quotient.  The only
division is by the global edge density $p$, which is at least $1/2$.
At $p=1$ the argument does not substitute into $1/(1-p)$; that endpoint
is handled separately.

The boundary of the theorem is exact: all non-spine vertices must be
independent page vertices adjacent to both ends of one common edge,
and every added vertex must be a private leaf attached to a page
vertex.  The proof does not attach leaves to arbitrary vertices of an
arbitrary positive graph.  It does not cover mixtures of page leaves
and spine leaves, edges among page vertices or leaves, shared leaves,
longer pendant paths, or books whose pages have more than one internal
vertex.  Such graphs require additional joint rooted estimates and are
not promoted by this result.

%%%%%%%%%%%%%%%%%%%%%%%%%%%%%%%%%%%%%%%%%%%%%%%%%%%%%%%%%%%%%%%%%%%%%%%%%%%%%%%
\appendix
\section{What the Lean formalization actually proves}
%%%%%%%%%%%%%%%%%%%%%%%%%%%%%%%%%%%%%%%%%%%%%%%%%%%%%%%%%%%%%%%%%%%%%%%%%%%%%%%

Corollary~\ref{cor:atlas} and Corollary~\ref{cor:atlas-cone} are formalized:
Atlas $41$, $114$, $138$ and $193$ are \texttt{verified}.  The development is
\texttt{Methods/Link/PageOp.lean} with the three rows in
\texttt{Methods/PageBook/}.  One difference of substance.

\subsection*{Lemma~\ref{lem:holder-compression} is not used, at any $m$}

The compression is the only place the argument leaves elementary territory,
and generalized H\"older does not appear in the development.  Every instance
the catalogue needs follows instead from the weighted Cauchy--Schwarz
inequality $(\int A\eta)^2\leq(\int A)(\int A\eta^2)$, taken at
$A(z)=W(x,z)W(y,z)$ and $\eta(z)=d(z)^{s/2}$, which the project already has:
\[
 H_{(a+b)/2}^2\leq H_aH_b\qquad(0\leq a\leq b),
\]
that is \texttt{sq\_pageOp\_le'}.  For the two-page rows $41$ and $114$ this is
\eqref{eq:holder-compression} at $m=2$, one application.  For the three-page
row $138$, where \eqref{eq:holder-compression} reads $H_{1/3}^3\leq H_0^2H_1$,
two applications suffice:
\[
 H_{1/3}^4\leq H_0^2H_{2/3}^2\leq H_0^2H_{1/3}H_1 ,
\]
from $H_{1/3}^2\leq H_0H_{2/3}$ and $H_{2/3}^2\leq H_{1/3}H_1$, and cancelling
one factor of $H_{1/3}$ (the case $H_{1/3}=0$ being immediate).  This is
\texttt{cube\_pageOp\_le}.

Lemma~\ref{lem:weighted-triangle} is formalized at real exponents, as the note
anticipates, and \texttt{weighted\_rootedTriangle\_rpow} sits alongside the
$\mathbb N$-indexed form that the clique common-leaf and cone families consume
unchanged.  \texttt{Real.rpow} sets $0^0=1$, which is the note's convention
$d^0=1$ including where the degree vanishes.

\end{document}
